\chapter{Umbilical Catheterization}

\section*{Overview}
Umbilical catheterization is the placement of a catheter into the umbilical artery or vein of a newborn for vascular access. It provides reliable arterial blood pressure monitoring and a route for blood sampling and the delivery of fluids and medications in critically ill neonates.

\section*{Alternative Names}
Umbilical arterial catheter; UAC; umbilical venous catheter; UVC; umbilical line

\section*{Principle}
The umbilical vessels remain patent in the newborn and provide direct access to the central circulation. The catheter tip is positioned in the descending aorta (arterial) or inferior vena cava (venous) under sterile technique, with tip location confirmed radiographically before use.

\section*{History}
Umbilical catheterization was introduced in the 1940s to 1960s and became an essential component of neonatal intensive care, enabling continuous arterial monitoring and safe administration of therapies to preterm and critically ill newborns.

\section*{Why This Test or Procedure Is Ordered}
Ordered in sick newborns requiring continuous arterial blood pressure monitoring, frequent arterial blood gas and laboratory sampling, exchange transfusion, parenteral nutrition, or the infusion of vasoactive medications and fluids when peripheral access is inadequate.

\section*{Is This a Primary or Secondary Test or Procedure}
A primary vascular access procedure in critically ill neonates and a secondary supportive measure when peripheral access is not feasible.

\section*{How This Test or Procedure Is Performed}
The umbilical stump is cleaned and draped sterilely, and the catheter is inserted into the umbilical artery or vein using sterile technique. The catheter is advanced to the predetermined depth, secured in place, and its position confirmed by radiography before infusion.

\section*{What Preparation Is Needed}
Informed consent from the parents or guardians, sterile preparation of the umbilical stump, and determination of appropriate insertion depth based on birth weight or shoulder-to-umbilicus length are needed before placement.

\section*{Contraindications and Precautions}
Necrotizing enterocolitis, omphalitis or infection of the umbilical stump, peritonitis, and lower-extremity vascular compromise are contraindications. Precautions include strict sterile technique, confirming catheter tip position, and discontinuing the line as soon as possible to reduce thrombotic and infectious risk.

\section*{Risks}
Risks include thrombosis or embolism, catheter-related bloodstream infection, vessel perforation with hemorrhage, air embolism, and extremity ischemia.

\section*{What Happens After the Test or Procedure}
The catheter is used for monitoring and infusion while in place, with limb perfusion, line position, and signs of infection monitored regularly. It is removed as soon as the infant's condition permits, typically after a few days.

\section*{Understanding Your Results}
Arterial catheter readings and samples guide respiratory and hemodynamic management, with results such as blood gases, glucose, and blood pressure interpreted against gestational-age-specific norms. Abnormal values are corrected by ventilation, fluids, and medications delivered through the line.

\section*{Normal Results}
Blood pressure, blood gas, glucose, and electrolyte values within the expected ranges for the infant's gestational and postnatal age indicate stable neonatal physiology.

\section*{Follow-up Tests and Procedures}
Peripheral intravenous or enteral access replaces the catheter once the infant is stable, and the site is observed for bleeding or infection after removal. Ongoing surveillance for complications such as hypertension or limb ischemia is warranted in selected cases.

\section*{References}
\begin{enumerate}
  \item MedlinePlus. Umbilical catheters. U.S. National Library of Medicine; 2025.
  \item Current Medical Diagnosis and Treatment. McGraw-Hill; 2025.
\end{enumerate}

\clearpage
